\documentclass{article}
\usepackage{graphicx} % Required for inserting images
\usepackage{amsmath}
\usepackage{amsfonts}
\usepackage{xcolor}
\usepackage{algorithmic}
\usepackage[lined,linesnumbered,ruled,vlined]{algorithm2e}


\title{FedHera: Towards Drift-Resilient Federated Fine-tuning with Heterogeneous Resources}
\author{Ke XIAO}
\date{November 2025}

\begin{document}

\maketitle
\section{Preliminaries and Motivation}
\subsection{Low-Rank Adaptation (LoRA) in FL}
There is currently a substantial body of research on federated fine-tuning methods based on LoRA. LoRA operates by adding trainable parameters to frozen pre-trained parameters, specifically a product of two low-rank matrices B and A. This low-rank property ensures the method's efficiency and resource-friendliness. Let $W_{0}$ denote the parameters of a pre-trained original model. LoRA aims to fine-tune these parameters on a target dataset, ultimately yielding fine-tuned model parameters $W$, which can be expressed as:
\begin{align}
    \mathbf{W} = \mathbf{W}_{0} + \Delta \mathbf{W} = \mathbf{W}_{0} + \mathbf{B}\mathbf{A},
\end{align}
where $\mathbf{W}_{0}, \Delta \mathbf{W} \in \mathbb{R}^{d \times p}$denote the pre-trained parameters and fine-tuned parameters respectively. The fine-tuning parameters can be expressed as the product of two low-rank matrices $\mathbf{B}$ and $\mathbf{A}$ where $\mathbf{B} \in \mathbb{R}^{d \times r} $ and $ \mathbf{A} \in \mathbb{R}^{r \times p}$ with $r \ll \min (d, p)$. During fine-tuning, the client maintains the pre-trained parameters frozen, and trains only the low-rank matrices $\mathbf{B}$ and $\mathbf{A}$ on local data to complete the process. This fine-tuning approach, which trains only two low-rank matrices while preserving the pre-trained parameters, offers considerable flexibility and with $r$ being far smaller than the pre-trained parameter dimension, LoRA proves highly accommodating for clients with limited computational resources.
\subsection{Heterogeneity and Truncation-based Methods}
Federal fine-tuning leverages precisely this advantage of LoRA, utilising data distributed across different clients to achieve fine-tuning of large models. Directly applying traditional FL to LoRA introduces challenges in model aggregation and addressing client heterogeneity. In traditional FL approaches such as FedAvg, the LoRA server aggregates models uploaded by clients through:
\begin{align}
    \mathbf{B}_{g} = \sum_{i \in \mathcal{N}}{\frac{1}{|\mathcal{D}_{i}|} \mathbf{B}_{i}},\quad 
    \mathbf{A}_{g} = \sum_{i \in \mathcal{N}}{\frac{1}{|\mathcal{D}_{i}|} \mathbf{A}_{i}}
\end{align}
However, due to the addition of the product of $\mathbf{B}$ and $\mathbf{A}$ to the pre-trained parameters in LoRA, this method of directly performing weighted averaging on $\mathbf{B}_{i}$ and $\mathbf{A}_{i}$ introduces a deviation from the ideal federated fine-tuning aggregation outcome:
\begin{align}
    \Delta \mathbf{W} = \sum_{i \in \mathcal{N}}{\frac{1}{|\mathcal{D}_{i}|} \mathbf{B}_{i}\mathbf{A}_{i}} \neq \mathbf{B}_{g}\mathbf{A}_{g}.
\end{align}
Additionally, traditional FL requires that $\mathbf{B}_{i}$ and $\mathbf{A}_{i}$ share identical dimensions to perform weighted averaging separately. This means all clients must accommodate the client with the least computational resources by selecting the same rank as it.

A large number of truncation-based researches have been proposed to address these issues. For instance, FlexLoRA employs a method where $\mathbf{B}_{i}$ and $\mathbf{A}_{i}$ are first multiplied and then weighted averaged, followed by singular value decomposition (SVD) on the aggregated $\Delta \mathbf{W}$. Server then distribute different-rank approximations of $\Delta \mathbf{W}$ to clients of varying tiers based on the matrix obtained from SVD. Specifically, FlexLoRA assigns distinct LoRA tiers ($r=\{8, 30, 200\}$) to resource-constrained clients, combining SVD to mitigate traditional FL aggregation bias while accommodating client heterogeneity.

The server first performs matrix multiplication on $\mathbf{B}_{i}$ and $\mathbf{A}_{i}$ from all clients, then computes the weighted sum to obtain the global update parameter $\mathbf{W}_{g}$. Immediately afterward, it performs SVD on it:
\begin{align}
    \mathbf{W}_{g} = \sum_{i \in \mathcal{N}}{\frac{1}{|\mathcal{D}_{i}|} \mathbf{B}_{i}\mathbf{A}_{i}},\quad SVD (\mathbf{W}_{g}) \rightarrow \mathbf{U}\mathbf{\Sigma} \mathbf{V}^{T}.
\end{align}
For each client's rank $r_i$, the server truncates from $\mathbf{U}, \mathbf{\Sigma}, \mathbf{V}$ and constructs $\mathbf{B}_{i}$ and $\mathbf{A}_{i}$ for client $i$:
\begin{align}
    \mathbf{U}_{i} = \mathbf{U}[:, :r_i], \mathbf{\Sigma}_{i}=&\mathbf{\Sigma}[:r_i, :r_i], \mathbf{V}_{i}=\mathbf{V}[:, :r_i],\\
    \mathbf{B}_{i}=\mathbf{U}_{i}\mathbf{\Sigma}^{1/2},&\quad \mathbf{A}_{i}=\mathbf{\Sigma}^{1/2}{\mathbf{V}_{i}}^{T}.
\end{align}
The server sends the constructed $\mathbf{B}_{i}$ and $\mathbf{A}_{i}$ to client $i$, thereby enabling LoRA federated fine-tuning with different ranks.

\subsection{Challenges in Heterogeneous Federated Fine-tuning}
\textbf{The Rigidity of Coupled Resource Allocation.} Current truncated methods typically assume that a client's communication capacity and computational capability are strongly coupled and linearly related, meaning that if a client can download a rank of $r=8$, it must also train a rank of $r=8$. In practical scenarios, these two factors are often mismatched. For instance, certain edge devices may possess sufficient bandwidth to download larger $\mathbf{B}_{i}$ and $\mathbf{A}_{i}$ matrices, yet their local memory capacity is so limited that only minuscule $\mathbf{B}_{i}$ and $\mathbf{A}_{i}$ can be trained locally and vice versa. Forcing both to be equal leads to significant resource wastage, either idling bandwidth or computational power. Furthermore, existing methods often apply the same rank across all layers, setting a fixed rank for clients of varying tiers at every level. However, different layers contribute differently to model performance. Simply truncating all layers according to a uniform rank leads to the discarding of parameters from critical layers, while less important layers still consume resources. A fine-grained strategy is required that decouples communication and computation across dimensions while also being able to perceive layer-wise importance during truncation.

\textbf{Truncation-Induced Accumulated Gradient Descent Drift (AGDD).} When employing this coarse-grained truncation method, a cumulative drift of non-negligible magnitude is introduced mathematically. When clients initialise based on the truncated subspace $W_{r_i}$, the computed gradient $\nabla f(W_{i})$ constitutes a biased estimate of the global true gradient $\nabla f(W)$. This bias is not merely an error from a single update but cascades across the $K$ epochs of local training. As $K$ increases, the trajectory of local model optimisation progressively deviates from the globally optimal subspace. To address this, incorporating more global information as constraints during local training becomes essential, imposing heightened demands on resource utilisation and the LoRA framework.

\section{Method: FedHera}
\subsection{Overview: FedHera Pipeline}
The foundation for supporting heterogeneity lies in the server performing SVD on the aggregated $\mathbf{W}_{g}$ and allocating LoRA matrices of varying ranks to clients through truncation. FedHera adopts the same processing approach. However, as $\mathbf{W}_{g}$ serves as the unified basis for the entire federated fine-tuning system, FedHera proposes an entirely new perspective on how to leverage this basis. Firstly, FedHera decouples the client's communication resources from its computational resources. Unlike existing methods that uniformly employ a truncate-and-train-all approach, FedHera achieves LoRA rank allocation for complex scenarios involving heterogeneous client resources. Furthermore, FedHera employs a finer-grained layer-wise truncation scheme, optimising limited resources by allocating them to layer parameters containing more global information. This reduces the AGDD caused by truncation while enhancing the effectiveness of federated fine-tuning. Finally, the introduction of a prefix-gating mechanism at the client further diminishes AGDD by providing clients with more global shared prior information during training.

As a federated fine-tuning framework, FedHera encompasses server-side parameter aggregation, distribution, and client-side training and uploading. Specifically, in each round, the server collects the client-uploaded $\mathbf{B}_{i}$ and $\mathbf{A}_{i}$ and performs aggregation and SVD according to Equation 4 (the initial round employs the same method as LoRA to generate the initial values of $\mathbf{B}$ and $\mathbf{A}$). Before distributing the aggregated parameters to clients, the server needs to assign a pair of ranks to each client via a dual-decoupling algorithm (Section 4.2). This guides the construction of the LoRA matrix sent to each client. Upon receiving the LoRA matrix and the ranks from the server, clients generate prefix masks for local training (Section 4.3) and upload the final training results to the server.
\subsection{Resource-Aware Heterogeneous Rank Allocation}
Given the practical scenario of heterogeneous clients, conflating communication and computational resources would result in wasteful utilisation of one or the other. FedHera generates a pair of ranks for each client at every layer of the LoRA model: $\{r_{tot},r_{main}\},\ 0 \leq r_{main} \leq r_{tot}$. Where, the communication rank $r_{tot}$ accounts for bandwidth constraints, determining how much information the client downloads from the server. The computational rank $r_{main}$ accounts for computational power constraints, determining how many parameters the client trains locally. When individuals or organisations within the same LAN perform federated fine-tuning, traditional approaches solely consider the constrained computational resources on the client to determine the rank of the LoRA model. FedHera, however, decouples these two resources. While clients train LoRA parameters with rank $r_{main}$, the difference $r_{tot} - r_{main}$ represents the additional global information supported by communication resources compared to conventional methods.

When considering constrained resources, traditional truncation methods do not account for layers individually, instead fixing the same rank for LoRA matrics across all layers. FedHera proposes a more granular, layer-wise normalised marginal energy-aware rank allocation scheme, assigning ranks to different layers of each client's LoRA model. The server performs singular value decomposition such that the decomposed $\mathbf{\Sigma}$ maintains singular values in descending order. This operation ensures FedHera possesses a fixed, unified basis for the current round on the server. The energy score $Energ_{l,k}$ for the $k$th singular value $\sigma_{l,k}$ of the $l$th layer:
\begin{align}
    Energ_{l, k} = \frac{\sigma_{l,k}^{2}}{\sum_{j}{\sigma_{l,j}^{2}}}.
\end{align}
With the limit on download bandwidth budget, FedHera must transmit as much global information as possible to clients. Specifically, it determines the rank $r^{tot}_{i, l}$ of the LoRA to be downloaded by client $i$ at layer $l$, based on client $i$'s download bandwidth budget $B_{i}$. Increasing the rank of the $l$th layer in the LoRA model by 1 incurs an additional transmission cost to client $i$: $c_{i,l}^{down} = (d_{l, out} + d_{l, in}) \cdot b_{i}$, where $d_{l, out}, d_{l, in}$ denote the parameter dimensions of layer $l$, and $b_{i}$ represents the effective payload byte count per scalar parameter sent from the server to client $i$ under different downlink quantisation schemes. The layer-normalised energy score $Score_{l,k}^{tot}$ for transmitting the $k$th singular value of layer $l$ to client $i$:
\begin{align}
    Score_{l,k}^{tot} = \frac{Energ_{l, k}}{c_{i,l}^{down}}.
\end{align}
Based on this score, FedHera employs a greedy water-filling algorithm targeting near-optimal strategies for discrete concave returns, Alg.~\ref{alg:rtot_waterfill}. Specifically, at each iteration, it selects the dimension with the largest $Score_{l,k}$ across all global layers and increments $r^{tot}_{i, l}$ by 1 until bandwidth is exhausted. This resolves resource wastage when computational resources fall below bandwidth constraints, ensuring the downloaded basis incorporates as much global information as possible.

\begin{algorithm}[tb]
\caption{Water-Filling for Download-side Rank Allocation $r^{\mathrm{tot}}_{i,\ell}$} \label{alg:rtot_waterfill}
\KwIn{Client $i$'s downlink budget $B_i$; per-layer caps $R_\ell$; per-layer costs $c^{down}_{i,\ell}$; singular values $\sigma_{\ell,k}$}
\KwOut{$r^{\mathrm{tot}}_{i,\ell}$}

Initialize residual budget $\widehat{B} \gets B_i$, ranks $r^{\mathrm{tot}}_{i,\ell} \gets 0,\ \forall \ell$ \\
\While{$\widehat{B} \ge \min\limits_{\ell:\ r^{\mathrm{tot}}_{i,\ell} < R_\ell} c^{down}_{i,\ell}$}{
    $best\_score \gets -\infty$, $best\_\ell \gets$ null \\
    \ForEach{layer $\ell$ with $r^{\mathrm{tot}}_{i,\ell} < R_\ell$}{
        $k \gets r^{\mathrm{tot}}_{i,\ell} + 1$ \\
        $\textit{score} \gets Energ_{l,k}/c^{down}_{i,\ell}$ \tcp*{$Score^{\mathrm{tot}}_{\ell,k}$}
        \If{$\textit{score} > best\_score$}{
            $best\_score \gets \textit{score}$; \quad $best\_\ell \gets \ell$
        }
    }
    \If{$best\_\ell =$ null}{\textbf{break}}
    \If{$\widehat{B} < c^{down}_{i,best\_\ell}$}{\textbf{break}}
    $r^{\mathrm{tot}}_{i,best\_\ell} \gets r^{\mathrm{tot}}_{i,best\_\ell} + 1$ \\
    $\widehat{B} \gets \widehat{B} - c^{down}_{i,best\_\ell}$ \\
}
\Return $r^{\mathrm{tot}}_{i,\ell}$
\end{algorithm}

Considering computational resources, given that client $i$ has time and memory constraints of $T_{i}$ and $M_{i}$ respectively, FedHera will allocate $r^{main}_{i, l}$ accordingly. The marginal computational cost per column is characterised by two slopes calibrated per client and per layer: the computational millisecond overhead per column, denoted as the per-column time slope: $c_{i,l}^{time} = \Delta t /\Delta r$, and the memory byte overhead per column, denoted as the per-column memory slope: $c_{i,l}^{mem} = \Delta m /\Delta r$. Both metrics are profiled by selecting a smaller $r$ and running multiple steps on the client under a fixed batch size. The layer normalisation energy score $Score_{l,k}^{main}$ when client $i$ adds an extra training pass for the $k$th singular value in layer $l$:
\begin{align}
    Score_{l,k}^{main} = \frac{Energ_{l, k}}{\alpha c_{i,l}^{time}+\beta c_{i,l}^{mem}},
\end{align}
Where $\alpha$ and $\beta$ are adaptive weights balancing the two resources. FedHera employs the same water-filling algorithm to assign $r^{main}_{i, l}$ to each layer for every client, Alg.~\ref{alg:rmain_waterfill}. The algorithm adjusts $\alpha$ and $\beta$ through adaptive residual budget allocation, adhering to client resource constraints and the condition of $\ 0 \leq r_{main} \leq r_{tot}$, thereby deriving feasible values for $r^{main}_{i, l}$.

\begin{algorithm}[tb]
\caption{Water-Filling for Train-side Prefix Rank Allocation}
\label{alg:rmain_waterfill}
\KwIn{Client $i$'s time/memory budgets $T_i,M_i$; caps $r^{\mathrm{tot}}_{i,\ell}$; slopes $c^{\mathrm{time}}_{i,\ell},c^{\mathrm{mem}}_{i,\ell}$; normalized energies $Energ_{\ell,k}$}
\KwOut{$r^{\mathrm{main}}_{i,\ell}$}

Initialize $\widehat T \gets T_i$, $\widehat M \gets M_i$; \quad $r^{\mathrm{main}}_{i,\ell}\gets 0,\ \forall \ell$ \\
\While{$\exists\,\ell:\ r^{\mathrm{main}}_{i,\ell}<r^{\mathrm{tot}}_{i,\ell}\ \wedge\ c^{\mathrm{time}}_{i,\ell}\le \widehat T\ \wedge\ c^{\mathrm{mem}}_{i,\ell}\le \widehat M$}{
    $best\_score \gets -\infty$;\ $best\_\ell \gets$ null \\
    $\textit{invT} \gets 1/\max(\widehat T,\varepsilon)$; \quad $\textit{invM} \gets 1/\max(\widehat M,\varepsilon)$ \\
    $\alpha \gets \textit{invT}/(\textit{invT}+\textit{invM})$; \quad $\beta \gets \textit{invM}/(\textit{invT}+\textit{invM})$ \\
    \ForEach{layer $\ell$ with $r^{\mathrm{main}}_{i,\ell}<r^{\mathrm{tot}}_{i,\ell}$}{
        \If{$c^{\mathrm{time}}_{i,\ell}\le \widehat T$ \textbf{and} $c^{\mathrm{mem}}_{i,\ell}\le \widehat M$}{
            $k \gets r^{\mathrm{main}}_{i,\ell}+1$ \\
            $\textit{score} \gets \dfrac{Energ_{\ell,k}}{\alpha\,c^{\mathrm{time}}_{i,\ell}+\beta\,c^{\mathrm{mem}}_{i,\ell}}$ \tcp*{$Score^{\mathrm{main}}_{\ell,k}$}
            \If{$\textit{score} > best\_score$}{
                $best\_score \gets \textit{score}$;\quad $best\_\ell \gets \ell$
            }
        }
    }
    \If{$best\_\ell =$ null}{\textbf{break}}
    $r^{\mathrm{main}}_{i,best\_\ell} \gets r^{\mathrm{main}}_{i,best\_\ell}+1$ \\
    $\widehat T \gets \widehat T - c^{\mathrm{time}}_{i,best\_\ell}$;\quad
    $\widehat M \gets \widehat M - c^{\mathrm{mem}}_{i,best\_\ell}$ \\
}
\Return $r^{\mathrm{main}}_{i,\ell}$
\end{algorithm}

\subsection{Drift Mitigation via Prefix-Gating and Faithful Aggregation}
Having computed the rank allocation strategy $\{\mathbf{r}_{i, l}^{tot}, \mathbf{r}_{i, l}^{main}\}$ for this global round, the server constructs LoRA-friendly factors $\mathbf{B}_{i}$ and $\mathbf{A}_{i}$ for each layer based on the singular value decomposition $\mathbf{U},\mathbf{\Sigma},\mathbf{V}^{T}$ and distributes them along with the corresponding $\mathbf{r}_{i, l}^{main}$ to client $i$. Upon receiving $\mathbf{B}_{i}$, $\mathbf{A}_{i}$ and $\{r_{i, l}^{tot}, r_{i, l}^{main}\}$, the client is accordingly initiated into prefix-gated training mode.

Specifically, the client defines a prefix gating operator 
$$G_{i,l}=\mathrm{diag}(\underbrace{1,\dots,1}_{r^{\mathrm{main}}_{i,l}},\underbrace{0,\dots,0}_{r^{\mathrm{tot}}_{i,l}-r^{\mathrm{main}}_{i,l}})$$
for each layer. During forward propagation, the downloaded complete adapter (of rank $r_{i, l}^{tot}$) remains engaged in inference, such that the output $h$ of any layer is:
\begin{align}
    h = \mathbf{W}_{0}x+s\mathbf{B}_{i}^{(l)}\mathbf{A}_{i}^{(l)}x,
\end{align}
in which $s$ denotes the scaling factor of LoRA. Although the local computational budget only supports training $r_{i, l}^{main}$ prefix sequences, the full $\mathbf{r}_{i, l}^{tot}$-dimensional space remains active during forward propagation. This ensures local optimisation consistently operates within a coordinate system aligned with the latest global principal subspace, thereby preventing representational gaps caused by truncation. Backpropagation employs prefix gating via $G_{i,l}$ to confine updates to the prefix subspace. The optimiser accumulates gradients and updates parameters solely for the first $r_{i, l}^{main}$ columns of $\mathbf{B}_{i}$ and the first $r_{i, l}^{main}$ rows of $\mathbf{A}_{i}$. The trailing $r_{i, l}^{tail}=r_{i, l}^{tot}- r_{i, l}^{main}$ columns yield zero gradients and are thus fixed as shared global residual information.

This design reduces computational and memory overhead, but more importantly, it suppresses cumulative gradient descent drift. Traditional truncated heterogeneous training isolates clients within low-dimensional truncated subspaces for optimisation, making it difficult to maintain alignment with the global direction. In FedHera, the frozen tail column provides a steady-state background aligned with the global principal components, whilst prefix gating constrains updates to this aligned subspace, thereby reducing variance accumulation arising from directional inconsistencies across clients. With local training complete, clients upload the updated $\mathbf{B}_{i}$ and $\mathbf{A}_{i}$ respectively. Given their low-rank properties, this step preserves LoRA's communication efficiency. The server then performs faithful aggregation across layers. Subsequently, SVD is reapplied to $W_g^{(l)}$ to refresh the unified basis, completing the closed-loop process from basis alignment to decoupled training and lossless aggregation.


\end{document}